
\documentclass[journal]{IEEEtran}

% packages
\usepackage{lipsum}
\usepackage[normalem]{ulem}
\usepackage{hyperref}
\usepackage{cite}
% \usepackage{flushend}
\usepackage[inline]{enumitem}
\usepackage{array,multirow,graphicx}
\usepackage{float}
\usepackage{orcidlink}

%% footnote
\usepackage{footnote}
\makesavenoteenv{tabular}
\makesavenoteenv{table}
\makesavenoteenv{tabular*}
\makesavenoteenv{table*}
% \usepackage{tablefootnote}

%% graphics packages
\usepackage{tikz}
\usepackage{colortbl}
\usepackage{pgfplots}
\usepackage{graphicx}
\usepackage{multirow}
\usepackage{circuitikz}
\usepackage{dblfloatfix}
\usepackage[bottom]{footmisc}
\ifCLASSOPTIONcompsoc
    \usepackage[caption=false, font=normalsize, labelfont=sf, textfont=sf]{subfig}
\else
    \usepackage[caption=false, font=footnotesize]{subfig}
\fi

%% checkmark packages 
\usepackage{pifont}
\newcommand{\cmark}{\ding{51}}
\newcommand{\xmark}{\ding{55}}
% tracked changes for review: removed text struck through in orange, added text in orange
\newcommand{\del}[1]{\leavevmode{\color{orange}\sout{#1}}}
\newcommand{\add}[1]{\leavevmode{\color{orange}\,#1}}

%% math packages
\usepackage{amssymb}
\usepackage{amsmath}
\usepackage{mathtools}
\usepackage{booktabs}

% math 
\IEEEeqnarraydefcolsep{0}{\hoffset}
\newcommand*{\defeq}{\stackrel{\text{def}}{=}}

% definition setup
%\theoremstyle{definition}

\newtheorem{defn}{Definition}[section]

% correct bad hyphenation here
\hyphenation{op-tical net-works semi-conduc-tor}

% nomenclature
\input{nomenclature}

% document
\begin{document}

% title, authors and thanks
\title{Robust Harmonic Hosting Capacity \\ Considering Parameter Uncertainty}
\author{Tom~Van~Acker~\orcidlink{0000-0001-6283-1457}, %
        Amauri~G.~Martins-Britto,~\IEEEmembership{Senior Member,~IEEE}~\orcidlink{0000-0002-2691-9061}, %
        and~Hakan~Ergun,~\IEEEmembership{Senior Member,~IEEE}~\orcidlink{0000-0001-5171-1986}%
\thanks{Tom Van Acker is with the Scenarios, Markets and Adequacy Department, Elia Transmission Belgium, 1000 Brussels, Belgium (e-mail: tom.vanacker@elia.com).}% <-this % stops a space
\thanks{Amauri G. Martins-Britto and Hakan Ergun are with the Research Group ELECTA, Department of Electrical Engineering, KU Leuven, 3000 Leuven, Belgium, and also with Etch-EnergyVille, 3600 Genk, Belgium (e-mail: amauri.martinsbritto@kuleuven.be, hakan.ergun@kuleuven.be).}
\thanks{This work is supported by the Energy Transition Fund by FOD Economy, Belgium, project HARMONIC. The General Board on Energy of the FOD Economy is not liable for any use of the presented work and subsequent results.}
\thanks{\textit{(Corresponding author: Tom~Van~Acker)}}}

% paper headers
\markboth{Journal of \LaTeX\ Class Files,~Vol.~14, No.~8, August~2015}%
{Shell \MakeLowercase{\textit{et al.}}: Bare Demo of IEEEtran.cls for IEEE Journals}

% make the title area
\maketitle

% abstract
\begin{abstract}
The harmonic hosting capacity of a power system depends on the network harmonic impedance, which follows from the physical inputs of its cables and lines, e.g., conductor temperatures, manufacturing tolerances and earth resistivity, known only to lie within intervals. This paper presents a robust harmonic hosting capacity method for phase-balanced three-phase power systems that considers uncertain harmonic current angles, through chance constraints, and cable and line parameters in the worst case over the intervals of their physical inputs. The parameters are computed at every harmonic frequency from these inputs, and the zero sequence includes the windings and earthing of the transformers. For fairness criteria that fix the ratios between the allocations of the units and a constant risk factor, the robust counterpart is a convex semi-infinite second-order cone program, solved by a cutting-set method with sign-predicted corners and resonance seeds. For the IEEE 30-bus system with underground cables and overhead lines, the uncertainty of the network costs 14--23\,\% of the hosting capacity up to the 25th harmonic, and the worst cases lie at the least damped network, near resonance. Scaling the fundamental parameters with the harmonic order misplaces the resonances, and its allocation violates the limits of the physical model, also when made robust over factor intervals. Up to the 50th harmonic, the transformer earthing changes the zero-sequence hosting capacity by up to 42\,\%, an effect that fades at the higher orders. Every case solves within minutes.
\end{abstract}

% keywords
\begin{IEEEkeywords}
Harmonic distortion, robust, fairness, hosting capacity, mathematical optimization.
\end{IEEEkeywords}

% peer-review title
\IEEEpeerreviewmaketitle

\section{Introduction}
\label{sec:introduction}

\subsection{Background and Motivation}
\label{sec:background_and_motivation}

\IEEEPARstart{R}{eaching} a net-zero economy by 2050 requires extensive electrification of the energy system. During the past two decades, renewable generation, particularly wind and photovoltaic power, has grown steadily and is expected to accelerate towards 2050~\cite{international_energy_agency_world_2023}. At the same time, decarbonization targets require the electrification of transport, heating, and industrial processes, replacing fossil fuels with electricity. In Europe, studies project that the total electricity demand will triple by 2050~\cite{energyville_paths_2023}. Consequently, future power systems will be largely dominated by power-electronic interfaced resources on both the supply and demand sides. This transition fundamentally changes the harmonic behavior of power systems. Traditional grids, dominated by synchronous machines and mostly linear loads, exhibited limited harmonic distortion. Modern systems, however, increasingly rely on converters whose nonlinear characteristics generate a broad spectrum of harmonic currents. The aggregation of numerous distributed converters creates new challenges in maintaining harmonic voltage quality across transmission and distribution networks.

\begin{table*}[!t]
    \centering
    \caption{Literature review on harmonic hosting capacity, network harmonic impedance uncertainty, bounded-uncertainty harmonic analysis and robust optimization over network parameters, in order of discussion, enumerating the uncertain features~(a)--(e). The entries marked $\star$ are prior work of the present authors.}
    \label{tab:literature_review}
    \begin{tabular}{ c | c | c | c c | c | c c c c c | r r }
        \multirow{2}{*}{reference} & \multirow{2}{*}{objective$^{a}$} & \multirow{2}{*}{method$^{b}$} & \multicolumn{2}{c |}{network} & \multirow{2}{*}{uncertainty$^{c}$} & \multicolumn{5}{c |}{uncertain features} & \multicolumn{2}{c}{tractability$^{d}$} \\
                                            &       &       & model & phases &       & (a)  & (b)   & (c)   & (d)   & (e)   & nodes & minutes \\ \hline
        \cite{noauthor_iec61000-3-6_2008}   & hhc   & ae    & nodal & bal.  & an.   &\cmark &\xmark &\xmark &\xmark &\cmark & --    & --    \\
        \cite{santos_considerations_2015}   & hhc   & ae    & nodal & bal.  & an.   &\cmark &\cmark &\xmark &\xmark &\xmark & --    & --    \\
        \cite{sun_distribution_2022}        & hhc   & cc    & full  & bal.  & scen. &\xmark &\cmark &\xmark &\cmark &\cmark & 16    & 30    \\
        \cite{nakhodchi_impact_2023}        & hhc   & ae    & nodal & bal.  & mc    &\cmark &\xmark &\xmark &\xmark &\xmark & 64    & --    \\
        \cite{ismael_probabilistic_2019}    & hhc   & mh    & full  & bal.  & mc    &\cmark &\cmark &\cmark &\cmark &\cmark & --    & --    \\
        \cite{emini_study_2026}             & hhc   & pf    & full  & bal.  & mc    &\xmark &\xmark &\xmark &\cmark &\xmark & --    & --    \\ \arrayrulecolor{gray!20} \hline
        \cite{testa_probabilistic_2002}     & hpf   & pf    & full  & bal.  & mc    &\cmark &\xmark &\cmark &\cmark &\xmark & 13    & --    \\
        \cite{li_probabilistic_2018}        & hra   & rma   & full  & bal.  & mc    &\xmark &\xmark &\cmark &\cmark &\xmark & 14    & 4.4   \\
        \cite{li2020}                       & hpf   & pf    & full  & bal.  & mc    &\cmark &\xmark &\cmark &\xmark &\xmark & 14    & --    \\
        \cite{barakou_parameter_2023}       & hra   & fs    & full  & --    & scen. &\xmark &\xmark &\cmark &\cmark &\xmark & --    & --    \\ \arrayrulecolor{gray!20} \hline
        \cite{romero_frequency_2018}        & hra   & fs    & full  & bal.  & pos.  &\xmark &\xmark &\xmark &\cmark &\xmark & 14    & --    \\
        \cite{huang_newton_2017}            & hpf   & pf    & full  & bal.  & int.  &\cmark &\xmark &\xmark &\xmark &\xmark & 30    & --    \\
        \cite{desousa_interval_2024}        & hpf   & pf    & full  & bal.  & int.  &\cmark &\xmark &\xmark &\cmark &\xmark & 33    & $<$\,0.01 \\
        \cite{lin_extended_2017}            & hra   & ae    & full  & bal.  & rob.  &\xmark &\xmark &\cmark &\xmark &\xmark & 81    & $<$\,0.001 \\ \arrayrulecolor{gray!20} \hline
        \cite{christakou_voltage_2018}      & vc    & ro    & lin.  & bal.  & rob.  &\xmark &\xmark &\cmark &\xmark &\xmark & 34    & --    \\
        \cite{liu_feasible_2023}            & fr    & ro    & lin.  & unb.  & rob.  &\xmark &\xmark &\cmark &\xmark &\xmark & --    & 8     \\ \arrayrulecolor{black} \hline
        $\star$~\cite{vanacker_hosting_2026}& hhc   & soc   & full  & bal.  & --    &\xmark &\xmark &\xmark &\xmark &\xmark & 1880  & $<$\,6 \\
        $\star$~\cite{vanacker_shhc}        & hhc   & qcqp/soc & full & bal. & pce  &\cmark &\xmark &\xmark &\xmark &\cmark & 30    & $<$\,0.04 \\ \arrayrulecolor{black} \hline
        \multirow{2}{*}{This work}          & \multirow{2}{*}{hhc}   & \multirow{2}{*}{soc}   & \multirow{2}{*}{full}  & \multirow{2}{*}{bal.}  & an./rob. &\cmark &\xmark &\cmark &\xmark &\cmark & 30$^{*}$ & $<$\,3 \\ 
        &&&&& rob. & \xmark & \xmark & \cmark &\xmark & \xmark & 30$^{**}$ & $<$\,8 \\
    \end{tabular}

    \vspace{2pt}
    \parbox{\textwidth}{\scriptsize
    $^{a}$Objective: (hhc)~harmonic hosting capacity, (hpf)~harmonic power flow, (hra)~harmonic resonance or impedance analysis, (vc)~voltage control and (fr)~feasible region of distributed energy resources, both at the fundamental frequency.
    $^{b}$Mathematical methods: (ae)~analytical expression, (pf)~power flow, (fs)~frequency scan, (rma)~resonance mode analysis, (mh)~metaheuristic optimization, (cc)~chance-constrained optimal power flow, (ro)~robust counterpart, (soc)~second-order cone optimal power flow, (qcqp)~quadratically constrained quadratic program. Network model: (nodal)~equivalent at one node, (full)~full network, (lin.)~linearized full network. Phases: (bal.)~balanced, (unb.)~unbalanced.
    $^{c}$Treatment of uncertainty: (mc)~Monte Carlo simulation, (an.)~analytical or moment-based, (scen.)~scenario-based, (pos.)~possibility theory, (int.)~interval or affine arithmetic, (pce)~polynomial chaos expansion, (rob.)~robust, i.e., guaranteed over a set.
    $^{d}$Tractability expressed as the number of nodes and the solve time in minutes. $^{*}$Positive- and negative-sequence harmonics 5--25, 23 uncertain physical inputs, on 10 threads. $^{**}$Zero-sequence variants without recombination, harmonics 2--50, on 16 threads.}
\end{table*}

A key challenge arises from the stochastic nature of the harmonic impedance of the network. System operators describe cables and overhead lines by their sequence resistance, reactance and susceptance at the fundamental frequency, taken from datasheets or commissioning measurements, and harmonic studies extrapolate these values to higher frequencies with scaling rules. In reality, the harmonic impedance follows from the physical inputs of each cable and line, e.g., the resistivities and temperatures of the conductors and screens, the permittivity of the insulation, the dimensions within their manufacturing tolerances, the laying geometry and the earth. ``In general, most of the physical parameters of a cable such as the resistivity, the permittivity and $\tan{\delta}$ are frequency- and temperature-dependent, and such data are not even available from a cable manufacturer. The same is true for the earth resistivity and permittivity.''~\cite[\S\,2.5.2.6]{ametani_cable_2015}. A cable project adds its own variability: the conductor temperature follows the load, the earth resistivity varies by orders of magnitude along the route and over the seasons, and the installed cables depart from their design within the acceptance limits of the standards. The harmonic impedances near a parallel resonance are the most sensitive to these inputs, and the scaling rules miss their joint effect. As the penetration of the converter increases, system operators need transparent metrics to assess whether additional units can be connected without exceeding harmonic limits. Harmonic hosting capacity provides such a metric by defining the maximum allowable harmonic emissions while ensuring compliance with power-quality standards.  The harmonic hosting capacity must therefore hold for every network that the physical inputs allow. However, most existing methods rely on impedance data at the fundamental frequency, extrapolated towards higher harmonics with scaling rules, leading to overly conservative or insufficiently protective results.

\subsection{Literature Review}
\label{sec:literature_review}

A growing body of literature addresses the harmonic hosting capacity of a power system, the uncertainty of the network harmonic impedance, harmonic analysis under bounded rather than probabilistic uncertainty, and optimization with a guarantee over uncertain data. Table~\ref{tab:literature_review} summarizes the publications discussed below, characterized by their objective, mathematical method, network representation, treatment of uncertainty and tractability, and by which of the following features they treat as uncertain:
\begin{enumerate*}
    \item [(a)] unit harmonic injection,
    \item [(b)] background harmonic distortion,
    \item [(c)] network parameters,
    \item [(d)] network exploitation, and
    \item [(e)] harmonic limits.
\end{enumerate*}

\subsubsection{Harmonic hosting capacity}

The standard IEC~61000-3-6 allocates planning levels over the installations connected to a point of common coupling and aggregates their emissions with an indicative summation law~\cite{noauthor_iec61000-3-6_2008}. Santos et al.\ used the summation law to determine the harmonic hosting capacity behind a Norton equivalent with background distortion of unknown phase~\cite{santos_considerations_2015}. Sun and Harrison proposed the first optimization-based stochastic method, with percentile-based harmonic distortion limits as chance constraints within an ac optimal power flow~\cite{sun_distribution_2022}. Nakhodchi et al.\ showed for a fast charging station for electric vehicles that the estimated hosting capacity depends strongly on how the medium-voltage network and the remote loads are modeled~\cite{nakhodchi_impact_2023}. Ismael et al.\ designed a filter that maximizes the harmonic-constrained hosting capacity of a photovoltaic system, with a metaheuristic around a Monte Carlo simulation over, among others, the filter parameters~\cite{ismael_probabilistic_2019}. At transmission level, a Monte Carlo study over the system loading found the harmonic hosting capacity to be governed by resonance-sensitive impedance~\cite{emini_study_2026}. The present authors determined an exact lower bound for the equitable harmonic hosting capacity with a second-order cone formulation of the full network~\cite{vanacker_hosting_2026}, showed that underground cables change it substantially~\cite{vanacker_impact_2025}, and extended it to stochastic harmonic current angles with chance constraints~\cite{vanacker_shhc}. Except for the filter in~\cite{ismael_probabilistic_2019} all of these methods take the network parameters as known.

\subsubsection{Uncertainty of the network harmonic impedance}
The joint CIGRE and CIRED guide stresses that the harmonic impedance changes continuously with load and network configuration and limits the validity of network models to about the 20th harmonic~\cite{robert_guide_1997}. Testa et al.\ built a probabilistic model of the nodal impedances from a 5\,\% tolerance on the inductances and capacitances, and concluded that harmonic voltage assessments that ignore the impedance variability give unreliable results~\cite{testa_probabilistic_2002}. Li et al.\ propagated a 10\,\% variation of the line, transformer and capacitor parameters through resonance mode analysis~\cite{xu_harmonic_2005,huang_application_2007} by Monte Carlo simulation, and obtained probabilistic resonance frequency bands and modal impedances~\cite{li_probabilistic_2018}. The same group separated aleatory from epistemic sources~\cite{li2020}. For offshore wind, Barakou et al.\ found that network and plant contingencies dominate the harmonic amplification, that the choice of cable model is noticeable and that a manufacturing tolerance of 8\,\% on the cable capacitance has a minor effect~\cite{barakou_parameter_2023}. Babaev et al.\ estimated the cable harmonic impedance online with an uncertainty budget, as its theoretical derivation neglects the installation~\cite{babaev_online_2020}. These works sample or scan the network parameters, but none of them optimizes the hosting capacity over their uncertainty.

\subsubsection{Harmonic analysis under bounded uncertainty}

Where probability distributions are not available, uncertainty is described by bounds. Romero et al.\ introduced the harmonic load flow based on possibility theory~\cite{romero_harmonic_2011} and extended the frequency scan with it to the possibility of a parallel resonance~\cite{romero_frequency_2018}. Interval and affine arithmetic bound the power flow solution for interval data~\cite{vaccaro_affine_2010} and were carried to harmonics. Huang et al.\ proposed a Newton-Raphson interval harmonic power flow in affine form~\cite{huang_newton_2017} and de Sousa et al.\ an interval harmonic power flow for microgrids based on a second-order Taylor expansion~\cite{desousa_interval_2024}. In all of these, the bounded quantities are the injections, while the network admittance is fixed. The paper presented by Lin et al.\ is the exception, based on the extended Gersgorin theorem, they derived a parameter domain that is free of harmonic resonance, which is conservative by construction~\cite{lin_extended_2017}. These methods bound a response or a feasible domain for given injections, but do not determine the hosting capacity.

\subsubsection{Robust optimization}

Robust optimization requires a constraint to hold for every realization of uncertain data within a set, and replaces the uncertain program by its robust counterpart~\cite{ben-tal_robust_1999}. Where no tractable counterpart exists, cutting-set methods alternate between a master problem over finite scenario sets and a worst-case analysis~\cite{mutapcic_cutting-set_2009}. In power systems, robust optimization has mainly addressed the uncertainty of renewable generation and load~\cite{roald_power_2023}, e.g., in robust ac and chance-constrained optimal power flow and hosting capacity~\cite{louca_robust_2019,lee_robust_2021,lubin_robust_2016,wang_distributed_2016}. Few works place the network itself in the uncertainty set. Christakou et al.\ controlled the voltage of a distribution network with line resistances that are bounded but of unknown distribution, due to conductor temperature and incomplete asset data, and linearized the network with sensitivity coefficients~\cite{christakou_voltage_2018}. Liu and Ma determined the robust feasible region of distributed energy resources for interval line impedances in a linear unbalanced power flow~\cite{liu_feasible_2023}. All of these are linear or linearized models at the fundamental frequency, whereas the network parameters enter the harmonic hosting capacity through the inverse of the harmonic admittance matrix, which is far from concave near a parallel resonance.

\subsubsection{Research gaps}

No existing method simultaneously achieves:
\begin{enumerate*}
    \item [(i)] an optimization of the harmonic hosting capacity with a guarantee over the uncertainty of the network parameters, i.e., robustly rather than by sampling,
    \item [(ii)] network parameters derived from the physical construction of the cables and lines and their frequency dependence, rather than from scaled fundamental-frequency data, and
    \item [(iii)] the joint treatment of both uncertainties, the angles in probability and the network in the worst case.
\end{enumerate*}
Moreover, the zero-sequence harmonics, for which the transformer windings and earthing set the network, are absent from the optimization-based methods.

% Table II (Section II) is defined here so that, as a double-column float, it lands at the bottom of the page on which Section II starts
\begin{table*}[!b]
    \centering
    \caption{Cables, overhead lines and earth of the IEEE 30-bus case: layout, and nominal value and interval of the 11 most influential of the 23 uncertain inputs, with the largest relative change of a sequence parameter at \textcolor{black}{50, 650 or 1250\,Hz} when the input alone moves to an end of its interval; $z_{0}$, $z_{1}$ and $y_{0}$, $y_{1}$ are the zero- and positive-sequence series impedance and shunt admittance per unit length. }
    \label{tab:parameter_box}
    \footnotesize\setlength{\tabcolsep}{4pt}
    \begin{tabular}{ >{\raggedright\arraybackslash}p{0.42\textwidth} >{\raggedright\arraybackslash}p{0.20\textwidth} l r }
        quantity & nominal & interval & largest change \\ \hline
        \multicolumn{4}{l}{\textit{cables, 33\,kV, 22 links in touching trefoil, center 0.8\,m deep, screens bonded at both ends}} \\
        aluminum conductor cross section~\cite{iec60228_2023} & 400--800\,mm$^{2}$ & by link & -- \\
        conductor DC resistance at 20\,$^{\circ}$C~\cite{iec60228_2023} & 0.0367--0.0778\,$\Omega$/km & $-8$\,\% to nominal & 4.4\,\% $z_{1}$ \\
        conductor diameter~\cite{iec60228_2023} & 23.0--34.2\,mm & nominal to $+0.4$\,mm & 1.3\,\% $y_{0}$ \\
        XLPE insulation thickness~\cite{iec60502_2_2014} & 8.0\,mm & 7.6 to 8.4\,mm & 4.4\,\% $y_{0}$ \\
        XLPE relative permittivity~\cite{iec60287_1_1_2006} & 2.4 & 2.1 to 2.5 & 12.5\,\% $y_{1}$ \\
        screen DC resistance at 20\,$^{\circ}$C, 48 copper wires of 1.35\,mm & 0.263\,$\Omega$/km & 0.251 to 0.263\,$\Omega$/km & 3.8\,\% $z_{0}$ \\
        jacket thickness~\cite{iec60502_2_2014} & 2.85--3.30\,mm & nominal to $+35$\,\% & 2.0\,\% $z_{1}$ \\
        conductor temperature~\cite{iec60502_2_2014} & 20\,$^{\circ}$C & 10 to 90\,$^{\circ}$C & 26\,\% $z_{0}$ \\
        \multicolumn{4}{l}{\textit{overhead lines, 132\,kV, 12 single-circuit lines in delta, span 300\,m}} \\
        ACSR phase conductors, 30 aluminum and 7 steel wires~\cite{en50182_2001} & Lion, Goat, Bear, Panther & by line & -- \\
        phase positions $(x, y)$ at the tower & $(\pm 2.5, 18.0)$, $(0, 22.33)$\,m & -- & -- \\
        earth wire, 7 steel wires of 3.15\,mm~\cite{iec63248_2022}, position & $(0, 29.0)$\,m & -- & -- \\
        reference sag at 20\,$^{\circ}$C, phases and earth wire & 6.0 and 5.4\,m & -- & -- \\
        phase temperature & 20\,$^{\circ}$C & $-5$ to 75\,$^{\circ}$C & 6.9\,\% $z_{1}$ \\
        ambient and earth-wire temperature & 20\,$^{\circ}$C & $-5$ to 45\,$^{\circ}$C & 3.2\,\% $z_{1}$ \\
        earth-wire steel relative permeability & 100 & 50 to 300 & 3.1\,\% $z_{0}$ \\
        \multicolumn{4}{l}{\textit{earth, homogeneous}} \\
        resistivity, logarithmic scale & 100\,$\Omega\cdot$m & 10 to 1000\,$\Omega\cdot$m & 14.6\,\% $z_{0}$ \\
        relative permittivity & 10 & -- & -- \\
        \multicolumn{4}{l}{\textit{12 further inputs}} \\
        semiconducting screens; wire diameters and lay ratios~\cite{en50182_2001,iec62641_2022,iec63248_2022}; metal resistivities; phase steel permeability & & & \textcolor{black}{$\leq 0.65$\,\%}
    \end{tabular}
\end{table*}

\subsection{Contributions and Outline}
\label{sec:contribution_and_outline}

In this paper, a holistic one-shot methodology is presented to determine the steady-state harmonic hosting capacity for an entire phase-balanced three-phase power system, extending~\cite{vanacker_shhc}. The presented method considers uncertainty on both: 
\begin{enumerate}
    \item [(i)] the angle of the injected harmonic current, enabling harmonic recombination, and
    \item [(ii)] cable and overhead-line parameters, 
\end{enumerate}
limiting the resulting individual and total harmonic voltage distortion, as well as the root-mean-square voltages and currents, using chance constraints. The problem is formulated from the perspective of the power system operator, who may anticipate the converter technologies its customers will connect, but not the operating point, parameters and control of the individual units. The angle of the harmonic current is therefore modeled as a random variable per unit. Additionally, the cable and overhead-line parameters follow from physical inputs, e.g., conductor temperatures, manufacturing tolerances, and earth resistances, known to lie within intervals. The harmonic hosting capacity is determined in terms of harmonic current magnitude using a fairness principle, anticipating angle uncertainty, and robust against parameter uncertainty. The robust counterpart over the parameter set is solved by a cutting-set method which gives a robust upper bound on the harmonic hosting capacity. The presented formulation uses the rectangular current-voltage variable space and incorporates:
\begin{itemize}
    \item uncertainty on the angle of the injected harmonic current via the angle model of~\cite{vanacker_shhc}, with exact moments in place of its truncated polynomial chaos expansion,
    \item series impedance and shunt admittance of the cables and overhead lines computed at every harmonic frequency from their physical inputs, with intervals from the acceptance limits of the standards and the operating ranges,
    \item exact representation of the harmonic impedance in the frequency domain and detailed treatment of transformers, in the zero, positive and negative sequence,
    \item enforcement of stochastic harmonic limits on bus voltage and component ampacity through chance constraints, and
    \item equitable distribution of the harmonic hosting capacity following a fairness principle.
\end{itemize}
For fairness criteria that fix the ratios between the allocations of the units, e.g., absolute equality, and a constant risk factor, the robust counterpart is a convex semi-infinite second-order cone program with one pair of variables per harmonic. The models and numerical illustrations are made freely available as part of the \textsc{Julia}-package \textsc{HarmonicPowerModels.jl}\footnote{Available at \url{https://github.com/timmyfaraday/HarmonicPowerModels.jl}.}, with the cable and line parameters from \textsc{LineCableModels.jl}\footnote{Available at \url{https://github.com/Electa-Git/LineCableModelsm.jl}.}~\cite{lcm_2026}.

The remainder of the paper is organized as follows: Section~\ref{sec:parameter_uncertainty} describes the physical inputs of the cables and overhead lines and their intervals. Section~\ref{sec:mathematical_framework} presents the robust counterpart of the stochastic harmonic hosting capacity problem and its solution by a cutting-set method. Section~\ref{sec:numerical_illustration} compares the risk factors and applies the method to the IEEE 30-bus system with physical cable and line inputs, including the zero-sequence harmonics and the transformer earthing. Section~\ref{sec:conclusion} concludes the paper.

\section{Cable and Overhead-Line Parameter Uncertainty}
\label{sec:parameter_uncertainty}

% $\mathcal{P}$ and $\bar{p}$ are \ParamSet and \ParamVec of the mathematical framework; use the macros once they move to nomenclature.tex.

The per-unit-length series impedance~$\bar{z}'_{\Line}(\Frequency;\bar{p})$ and shunt admittance~$\bar{y}'_{\Line}(\Frequency;\bar{p})$ follow from physical inputs of the cables and overhead lines, accounting for the corresponding material and dimensional tolerances, and computed with the Julia package \textsc{LineCableModels.jl}~\cite{lcm_2026}. Each cable or line is a system of tubular conductors separated by insulating and semiconducting layers, in the air or in the earth~\cite{ametani_cable_2015}. For~$n$ conductors, the $n \times n$ matrices of series impedance and shunt admittance per unit length are
\begin{IEEEeqnarray}{ r C l }
    Z'(\Frequency;\bar{p}) &=& Z'_{\text{c}} + Z'_{\text{i}} + Z'_{\text{e}}, \label{eq:pul_impedance} \\
    Y'(\Frequency;\bar{p}) &=& \mathrm{j} 2\pi\Frequency \left( P'_{\text{i}} + P'_{\text{e}} \right)^{-1}, \label{eq:pul_admittance}
\end{IEEEeqnarray}
in which $Z'_{\text{c}}$ collects the surface impedances of the conductors, $Z'_{\text{i}}$ the impedances of the insulating layers and $Z'_{\text{e}}$ the earth-return impedances; and $P'_{\text{i}}$ and $P'_{\text{e}}$ are the potential coefficients of the insulating layers and of the surrounding air and earth. \textcolor{black}{These expressions hold from the fundamental to frequencies well above 2.5\,kHz, the 50th harmonic~\cite{ametani_cable_2015}. Specifically, the formulations adopted in LineCablemodels.jl have been validated against the finite element method (FEM) up to 100 MHz.}

% The surface impedances are ratios of modified Bessel functions of~$r\sqrt{\mathrm{j}2\pi\Frequency\mu/\rho}$ for a conductor of radius~$r$, permeability~$\mu$ and resistivity~$\rho$. The factor~$\sqrt{\pi\Frequency\mu/\rho}$ is the reciprocal of the skin depth, which in aluminum falls from 12.0\,mm at 50\,Hz to 3.3\,mm at 650\,Hz, below the 13.25\,mm radius of a 500\,mm$^{2}$ core. The earth return penetrates 712\,m at 50\,Hz and 197\,m at 650\,Hz into a 100\,$\Omega\cdot$m earth, so the zero-sequence parameters, whose currents return through screens, earth wires and earth, vary most. 
The effect of temperature~$\vartheta$, in $^{\circ}$C, on metal resistivity is described by 
\begin{equation}
    \rho(\vartheta) = \rho_{20} \left[ 1 + \kappa_{20} \left( \vartheta - 20 \right) \right],
    \label{eq:resistivity_temperature}
\end{equation}
where $\rho_{20}$ is the resistivity at 20\,$^{\circ}$C and $\kappa_{20}$ its temperature coefficient~\cite{iec62641_2022,iec60287_1_1_2006}. In overhead lines, temperature also sets the sag through the elastic thermal change-of-state equation, and each conductor lies at its support height less two thirds of the sag.

Cable screens are assumed bonded and solidly earthed at both cable terminals and the earth wires at every tower, with circuits ideally transposed. This yields circulant-reduced $3 \times 3$ matrices, expressed as zero-, positive- and negative-sequence parameters ~$\bar{z}'_{\Line}$ and~$\bar{y}'_{\Line}$ of Section~\ref{sec:parametric_network}.

Table~\ref{tab:parameter_box} summarizes the cables, lines and earth of the IEEE 30-bus case, whose cable sizes and ACSR codes follow from the branch currents, rated under IEEE~738~\cite{ieee738_2023}. The intervals are the acceptance limits of the standards cited in Table~\ref{tab:parameter_box}, operating and environmental ranges for the temperatures and the earth, and study intervals otherwise. Two continuous maps, equal to the identity at feasible points, keep every point of the box physical: the phase temperature is at least the ambient one, and the wire-diameter deviations of each ACSR code are scaled back onto the diameter and mass limits of EN~50182~\cite{en50182_2001}.

Physical inputs are represented as parameter space coordinates per input of a family, e.g., one conductor temperature for all 22 cables or, for a per-size quantity, it moves every size within its own interval. Of 39 candidate inputs, the 23 that change some sequence impedance or admittance by at least 0.1\,\% at \textcolor{black}{50, 650 or 1250\,Hz}, when moved alone to an end of their interval, form the \textcolor{black}{box coordinates}. The complete dataset, with the source and class of every value, and the specifications and case files are available in the HarmonicPowerModels.jl repository~\cite{hpm_2026}.

\input{journals/rhhc/sections/mathematical_framework}
\section{Numerical Illustration}
\label{sec:numerical_illustration}

This section presents the numerical illustration of the robust harmonic hosting capacity problem of a phase-balanced three-phase power system. Two case studies are considered on the same test network: first, every harmonic order up to the 50th, including the zero-sequence harmonics, without uncertainty on the current angles, to assess the transformer windings and earthing; and second, the characteristic harmonics up to the 25th with uncertain current angles, to assess the choice of the risk factor and the network data and their uncertainty. All results are presented in per-unit using a 100\,MVA basis.

\pgfplotsset{rhhc/.style={
    axis lines=left, xmajorgrids=true, ymajorgrids=true, grid style={gray!10},
    label style={font=\footnotesize}, tick label style={font=\footnotesize},
    every axis plot/.append style={line width=0.6pt, mark size=1.5pt, black},
    legend style={draw=none, font=\scriptsize, legend columns=-1, at={(0.5,1.02)}, anchor=south},
    legend cell align=left},
    reference/.style={gray!60, line width=1.2pt, mark=none}}

\subsection{Test Network}
\label{sec:test_network}

The test network is the IEEE 30-bus system of~\cite{vanacker_shhc}, with 21 loads as harmonic units under the absolute equality fairness criterion. The per-unit-length parameters of Section~\ref{sec:parameter_uncertainty} are computed with \textsc{LineCableModels.jl}~\cite{lcm_2026} in the exact $\pi$-model~\eqref{eq:pi_series}--\eqref{eq:pi_shunt}. The 33\,kV lines are replaced by 22 underground cables of four sizes, 400 to 800\,mm$^{2}$ aluminum, 18/30(36)\,kV, in touching trefoil at 0.8\,m depth. The twelve 132\,kV lines are aluminum conductors steel reinforced of four sizes on a delta tower under one steel earth wire. Cable screens are solidly bonded at both ends, the earth wire is grounded at every tower, and all circuits are ideally transposed. The transformers carry their vector groups, i.e., YNyn0 for the 132/33\,kV units and YNd11 for the condenser step-up units. The fundamental operating point is that of the nominal parameters. Four network models are compared:
\begin{enumerate}
    \item[(i)] \textit{placeholder}, a single 630\,mm$^{2}$ aluminum cable design for all 33\,kV links, with four uncertain inputs, and the 132\,kV lines of~\cite{vanacker_shhc} without line charging;
    \item[(ii)] \textit{cables}, the four cable sizes with four uncertain inputs, i.e., resistance, insulation thickness, permittivity and earth resistivity, and the same 132\,kV lines;
    \item[(iii)] \textit{cables and lines}, case~(ii) with the 132\,kV lines modeled physically, including their charging; and
    \item[(iv)] \textit{full}, case~(iii) with all 23 physical inputs of Section~\ref{sec:parameter_uncertainty}.
\end{enumerate}

\subsection{Zero-Sequence Harmonics and Transformer Earthing}
\label{sec:zero_sequence_results}

The first case study considers network model~(iii) with every harmonic order from 2 to 50, including the 16 zero-sequence orders, without recombination, as in~\cite{vanacker_hosting_2026}: the unit contributions align in the worst case, so that no angle model of the units and no risk factor are needed. The 50th harmonic, at 2.5\,kHz, lies within the valid range of the cable and line model of Section~\ref{sec:parameter_uncertainty}. Each generator is taken to connect through a YNd step-up with a leakage reactance of 0.12\,pu on a rating of 0.2 times its short-circuit power, at~$X/R = 40$. The synchronous condensers at buses~11 and~13 sit behind their own step-up units, whose delta windings block the zero sequence. % generator zero-sequence data: assumption accepted by Tom Van Acker on 2026-11-11
Five variants differ only in the windings and earthing of the transformers with a 33\,kV winding:
\begin{enumerate*}
    \item[(a)] YNyn0 with solidly earthed 33\,kV stars;
    \item[(b)] YNyn0 with the 33\,kV stars earthed through 19.05\,$\Omega$, which limits an earth fault to 1\,kA, i.e., case~(iii);
    \item[(c)] Dyn11, with the delta on the 132\,kV side and the 33\,kV stars as in~(b);
    \item[(d)] YNd11, with the delta on the 33\,kV side and the 33\,kV star of the condenser step-up unearthed; and
    \item[(e)] YNy0 with every 33\,kV star unearthed.
\end{enumerate*}
The clock numbers keep the same phase shift on every 132/33\,kV path.

\begin{figure}[b]
    \centering
    \begin{tikzpicture}
        \begin{semilogyaxis}[rhhc, scale only axis, width=0.78\columnwidth, height=2.5cm, xmin=1.5, xmax=50.5, ymin=3e-5, ymax=4e-3,
            xtick={2,5,10,15,20,25,30,35,40,45,50}, xlabel={harmonic order $\Harmonic$}, ylabel={$\HostingCapacity_{\Harmonic}$ [pu]},
            every axis plot/.append style={mark size=1.1pt, line width=0.5pt},
            xlabel style={yshift=3pt}, ylabel style={yshift=-3pt},
            legend image code/.code={\draw[mark repeat=2, mark phase=2] plot coordinates {(0cm,0cm) (0.14cm,0cm) (0.28cm,0cm)};},
            legend columns=-1, legend style={at={(0.5,1.02)}, anchor=south, draw=none, inner sep=0pt, font=\scriptsize, /tikz/every even column/.append style={column sep=1.5pt}, /tikz/every odd column/.append style={column sep=-1pt}}]
            \pgfplotsinvokeforeach{3,6,...,48}{\draw[gray!18, line width=3.6pt] (axis cs:#1,3e-5) -- (axis cs:#1,4e-3);}
            \addplot[mark=o, solid]               table[x=h, y=ra] {journals/rhhc/data/zero_sequence.dat};
            \addplot[mark=square*, densely dashed] table[x=h, y=rb] {journals/rhhc/data/zero_sequence.dat};
            \addplot[mark=triangle*, dotted]       table[x=h, y=rc] {journals/rhhc/data/zero_sequence.dat};
            \addplot[mark=diamond*, dashdotted]    table[x=h, y=rd] {journals/rhhc/data/zero_sequence.dat};
            \addplot[mark=x, densely dotted, mark size=1.6pt] table[x=h, y=re] {journals/rhhc/data/zero_sequence.dat};
            \legend{(a)\,14.3, (b)\,13.0, (c)\,13.5, (d)\,12.8, (e)\,12.8}
        \end{semilogyaxis}
    \end{tikzpicture}
    \vspace{-1.5ex}
    \caption{Robust harmonic hosting capacity per harmonic without recombination, earthing variants~(a)--(e), network model~(iii), $d = 4$; legend: total robust hosting capacity~$\sum_{\Harmonic \in \HarmonicSet} \HostingCapacity_{\Harmonic}$ in~$10^{-3}$\,pu; gray: zero-sequence harmonics.}
    \label{fig:zero_sequence}
\end{figure}

Fig.~\ref{fig:zero_sequence} shows that the earthing acts on the zero sequence only. At the positive- and negative-sequence harmonics, the robust hosting capacity of the five variants agrees to within~$10^{-6}$\,pu.

Solid earthing of the 33\,kV stars hosts~42\,\% more zero-sequence current than the earthing through 19.05\,$\Omega$ after robustness. Removing the earthing altogether, variants~(d) and~(e), lowers the robust zero-sequence hosting capacity by only a further 6--8\,\%, as the resistance already decouples the 33\,kV network. A delta on the 132\,kV side, variant~(c), raises it by 15\,\%, as the 132\,kV network no longer shares the zero-sequence path. The effect of the earthing fades with the harmonic order: up to the 24th harmonic, the robust zero-sequence hosting capacity differs by a factor of up to 5.7 between the variants, from the 33rd harmonic on by at most 12\,\%, as the zero-sequence current returns increasingly through the capacitances of the cables and lines rather than the transformer neutrals: per unit injected at a unit bus, the current through the transformer earthing falls from 0.39--0.68 at the 3rd harmonic to at most 0.07 at the 45th. The earthing practice of the distribution network is thus a lever on the zero-sequence hosting capacity of the order of the uncertainty of the network itself.

The uncertain cable parameters cost 9--16\,\% of the zero-sequence hosting capacity, against 6--15\,\% for the positive- and negative-sequence harmonics. The loss concentrates at a few high orders, the 24th and 30th harmonic, where the robust hosting capacity falls to 0.57--0.67 of the nominal one, depending on the earthing. The zero-sequence parameters of the cables and the earthing path together place the zero-sequence resonances, so a deterministic study of the zero-sequence harmonics with nominal cable data misses the harmonics that are most at risk.

\subsection{Harmonics with Uncertain Current Angles}
\label{sec:ieee30_physical}\label{sec:risk_factor_results}

The second case study considers the characteristic harmonics up to the 25th,  with recombination. The current angle of every unit is a random variable, affine in one Beta-distributed random variable per unit, with the angle distributions of~\cite{vanacker_shhc}, so that the angles are independent between the units and fully dependent across the harmonics of a unit. The unit contributions~\eqref{eq:response} recombine at every bus and branch, and the chance constraints~\eqref{eq:cc_compact}, with the exact moments~\eqref{eq:mean_vector}--\eqref{eq:gram_matrix}, limit the result at the risk level~$\varepsilon = 0.05$. The robust hosting capacity is computed for the four network models~(i)--(iv), with the 21 units as diode bridges. First, the choice of the risk factor is assessed on the purely inductive network of~\cite{vanacker_shhc}, without parameter uncertainty, for the diode bridges. 

The Gaussian and Cantelli factors (\S\,\ref{sec:shhc_recap}) are compared on the nominal purely inductive IEEE 30-bus system, i.e., without parameter uncertainty, where the exact~$(1-\varepsilon)$ quantile follows from Monte Carlo simulation. Both act on the exact moments. Fig.~\ref{fig:risk_factor} compares them with the exact quantile for a risk level~$\varepsilon \in \{0.01, 0.02, 0.05, 0.1, 0.2\}$. The Cantelli bound never exceeds the risk level, with violation probabilities of at most 5\,\% of~$\varepsilon$, but costs~2.6--13\,\% of hosting capacity. The Gaussian factor stays within~1.0\,\% of the exact quantile, with a violation probability below the risk level for~$\varepsilon \leq 0.1$ and of~1.07 times the risk level at~$\varepsilon = 0.2$. The Gaussian factor on exact moments is therefore used in the remainder of this section.

\begin{figure}[!b]
    \centering
    \pgfplotsset{rfplot/.style={rhhc, scale only axis, width=0.3\columnwidth, height=2.0cm, xmin=0.008, xmax=0.25,
        xtick={0.01,0.05,0.2}, xticklabels={0.01,0.05,0.2}, xlabel={risk level $\varepsilon$}},
        can/.style={solid, mark=o}, gau/.style={densely dashed, mark=square*}}
    \pgfplotslegendfromname{rflegend}\\[-2pt]
    \subfloat[hosting capacity relative to the exact $(1-\varepsilon)$ quantile]{%
    \begin{tikzpicture}
        \begin{semilogxaxis}[rfplot, ymin=-15, ymax=2, ylabel={difference [\%]},
            legend columns=3, legend to name=rflegend, legend style={font=\scriptsize, /tikz/every even column/.append style={column sep=4pt}}]
            \addplot[reference] coordinates {(0.008,0) (0.25,0)};
            \addplot[can] table[x=eps, y=ddcan] {journals/rhhc/data/risk_factor.dat};
            \addplot[gau] table[x=eps, y=ddgau] {journals/rhhc/data/risk_factor.dat};
            \legend{target, Cantelli, Gaussian}
        \end{semilogxaxis}
    \end{tikzpicture}}\hfill
    \subfloat[violation probability of the binding limit]{%
    \begin{tikzpicture}
        \begin{loglogaxis}[rfplot, ymin=1e-4, ymax=0.3, ylabel={probability [-]}, ytick={1e-3,1e-2,1e-1}, extra y ticks={1.5e-4}, extra y tick labels={0}, extra y tick style={grid=none}]
            \addplot[reference] coordinates {(0.008,0.008) (0.25,0.25)};
            \addplot[can] table[x=eps, y expr={max(\thisrow{vdcan},1.5e-4)}] {journals/rhhc/data/risk_factor.dat};
            \addplot[gau] table[x=eps, y=vdgau] {journals/rhhc/data/risk_factor.dat};
        \end{loglogaxis}
    \end{tikzpicture}}
    \caption{Choice of the risk factor, nominal IEEE 30-bus system, $N = 21$ diode bridges, harmonics up to the 25th: (a) total hosting capacity with the Cantelli and Gaussian factor on exact moments relative to the exact $(1-\varepsilon)$ quantile, and (b) violation probability of the binding individual harmonic voltage limit by Monte Carlo simulation with 20\,000 samples, against the risk level~$\varepsilon$ (gray: target, i.e., the exact quantile in~(a) and a violation probability equal to~$\varepsilon$ in~(b); violation probabilities of zero at the bottom).}
    \label{fig:risk_factor}
\end{figure}

Fig.~\ref{fig:ieee30_physical} gives the results for
\begin{enumerate*}
    \item [(a)] the total harmonic hosting capacity for the four network models, and
    \item [(b)] the per-harmonic robust hosting capacity relative to the nominal stochastic one.
\end{enumerate*}
No case shows a violation on 1024 random points of~$\ParamSet$. At the worst-case points of the final scenario sets of case~(iv), the sampled 95\,\% quantile of the binding limits lies between 0.5\,\% below and 3.0\,\% above the limit, and their violation probability between 0.035 and 0.064, i.e., at most 1.3 times the risk level.

\begin{figure}[!tb]
    \centering
    \pgfplotsset{nmplot/.style={rhhc, scale only axis, width=0.25\columnwidth, height=2.2cm,
        legend style={at={(0.5,1.03)}, anchor=south, draw=none, font=\scriptsize, /tikz/every even column/.append style={column sep=2pt}},
        legend image code/.code={\draw[mark repeat=2, mark phase=2] plot coordinates {(0cm,0cm) (0.15cm,0cm) (0.3cm,0cm)};}},
        mi/.style={solid, mark=o}, mii/.style={densely dashed, mark=square*}, miii/.style={dotted, mark=triangle*}, miv/.style={dashdotted, mark=diamond*},
        res/.style={only marks, mark=|, mark size=2.5pt, line width=0.7pt, forget plot}}
    \subfloat[total]{%
    \begin{tikzpicture}
        \begin{axis}[nmplot, ybar, bar width=4pt, ymin=0, ymax=0.165, ylabel={$\sum_{\Harmonic} \HostingCapacity_{\Harmonic}$ [pu]}, ylabel style={yshift=-4pt},
            symbolic x coords={(i),(ii),(iii),(iv)}, xtick=data, enlarge x limits=0.18, legend columns=2, xlabel={network model},
            legend image code/.code={\draw[#1] (0cm,-0.08cm) rectangle (0.2cm,0.08cm);},
            scaled y ticks=false, yticklabel style={/pgf/number format/fixed, /pgf/number format/precision=2},
            nodes near coords style={font=\tiny, rotate=90, anchor=west, inner sep=1pt}]
            \addplot[black, fill=white] coordinates {((i),0.1330) ((ii),0.1282) ((iii),0.0993) ((iv),0.0993)};
            \addplot[black, fill=gray!60, nodes near coords, point meta=explicit symbolic]
                coordinates {((i),0.1145)[0.86] ((ii),0.1038)[0.81] ((iii),0.0797)[0.80] ((iv),0.0769)[0.78]};
            \legend{nominal, robust}
        \end{axis}
    \end{tikzpicture}}\hfill
    \subfloat[ratio per harmonic]{%
    \begin{tikzpicture}
        \begin{axis}[nmplot, xmin=4, xmax=26, ymin=0.3, ymax=1.05, xtick={5,7,11,13,17,19,23,25}, xticklabels={5,,11,,17,,23,},
            xlabel={harmonic order $\Harmonic$}, ylabel={ratio [-]}, legend columns=4, ylabel style={yshift=-4pt},
            every axis plot/.append style={mark size=1.3pt}]
            \addplot[reference, forget plot] coordinates {(4,1) (26,1)};
            \addplot[mi]   coordinates {(5,0.834) (7,0.896) (11,0.865) (13,0.899) (17,0.861) (19,0.866) (23,0.845) (25,0.854)};
            \addplot[mii]  coordinates {(5,0.906) (7,0.966) (11,0.782) (13,0.851) (17,0.866) (19,0.761) (23,0.747) (25,0.758)};
            \addplot[miii] coordinates {(5,0.884) (7,0.483) (11,0.885) (13,0.623) (17,0.867) (19,0.989) (23,0.929) (25,0.738)};
            \addplot[miv]  coordinates {(5,0.825) (7,0.447) (11,0.874) (13,0.581) (17,0.859) (19,0.937) (23,0.868) (25,0.711)};
            \addplot[res, gray] coordinates {(6.86,0.36) (18.22,0.36) (6.95,0.36) (18.35,0.36)};
            \addplot[res, black] coordinates {(6.27,0.43) (12.32,0.43) (16.86,0.43) (18.21,0.43) (19.12,0.43) (19.65,0.43) (22.72,0.43)};
            \legend{(i), (ii), (iii), (iv)}
        \end{axis}
    \end{tikzpicture}}
    \caption{Nominal stochastic and robust harmonic hosting capacity of the IEEE 30-bus system for the network models~(i)--(iv), harmonics up to the 25th, Gaussian factor on exact moments, $\varepsilon = 0.05$: (a) total, with the robust-to-nominal ratio above the bars, and (b) robust-to-nominal ratio per harmonic, with the parallel resonances at the nominal parameters as ticks at the bottom, gray for~(i) and~(ii), black for~(iii) and~(iv).}
    \label{fig:ieee30_physical}
\end{figure}

With cables only, i.e., case~(ii), the parallel resonances lie at~$\Harmonic = 6.95$ and~18.4. Modeling the 132\,kV lines physically, i.e., case~(iii) moves the first resonance to~$\Harmonic = 6.27$ and adds resonances at~$\Harmonic = 12.32$ and between the 17th and the 23rd harmonic. The robust hosting capacity relative to the nominal one halves at the 7th harmonic, from 0.966 to 0.483, and drops from 0.851 to 0.623 at the 13th harmonic. The network data that enter the robust problem thus matter as much as their uncertainty. A simplified network places the resonances, and with them the worst cases, at other harmonics.

Extending the parameter set from 4 to 23 physical inputs, i.e., case~(iv), lowers the robust hosting capacity by a further~3.5\,\%, mostly at the 5th, 7th, 13th, 19th and 25th harmonic. The uncertainty of the inputs beyond the four dominant ones is therefore not negligible, but it is second order. With~$d = 23$, the dictionary holds 191 sign-predicted corners in place of the~$2^{23}$ vertices. Without them, the method signs off after five iterations at a solution that violates a limit by~0.16\,\% at a sign-predicted corner, whereas 1024 independent random points of~$\ParamSet$ find no violation. With the corners, no violation remains. In high dimension, random sampling, including the a-posteriori check, does not certify robustness, the corners do to first order. The worst corner of the binding limit, the 23rd harmonic at bus~3, is the least damped network, i.e., the coldest cables and lines and the lowest earth resistivity. Moreover, the robust-to-nominal ratio ranges from 0.45 to 0.99 over the harmonics, Fig.~\ref{fig:ieee30_physical}(b). A robust budget is therefore a new allocation, not a uniform derating of the nominal one.

A common simplification scales the fundamental $\pi$-model with the harmonic order, i.e., the resistance with~$\sqrt{\Harmonic}$ and the reactance and susceptance with~$\Harmonic$. Against the physical model, Fig.~\ref{fig:scaling}(a), the rule overestimates the resistance of the lines by a factor of two to three and the reactance of the cables by up to 2.4 at the 25th harmonic, as the screen return lowers their inductance, while the physical cable resistance exceeds the rule by up to 36\,\%. The scaled model thereby moves the parallel resonance near the 23rd harmonic, where the binding limit lies, and halves the driving-point impedance at bus~3 there, Fig.~\ref{fig:scaling}(b). Applied to the nominal 50\,Hz parameters of case~(iv), it lowers the nominal hosting capacity by~21\,\%, yet its allocation over the harmonics violates the limits of the physical model at the nominal parameters. Factor intervals derived from the same box, i.e., the hull of the ratios of the resistance, inductance and capacitance of each family to their nominal values over the box, at the fundamental and at every harmonic frequency, give a robust hosting capacity of 0.0482\,pu, 37\,\% below the physical one, which still violates the limits of the physical model at every one of 1024 random points of~$\ParamSet$. The same factors applied to the nominal frequency response of \textsc{LineCableModels.jl} give 0.0726\,pu, 5.7\,\% below the physical one, with no violation of the physical model at 1024 random points. The frequency dependence of the cable and line parameters is therefore part of the robust problem, and a box of factors on the 50\,Hz parameters does not contain the physical uncertainty. \add{In the zero sequence, the scaling rule fails altogether, as the resistance follows the return path rather than the skin effect of the conductor: up to the 15th harmonic, it rises by only 26\,\% for the cables, whose screens carry the return, and 6.5 times for the lines, whose return runs through the earth, against 3.9 times for~$\sqrt{\Harmonic}$.}

\begin{figure}[!tb]
    \centering
    \pgfplotsset{scplot/.style={rhhc, scale only axis, width=0.29\columnwidth, height=2.2cm, xlabel={harmonic order $\Harmonic$},
        legend columns=2, legend style={at={(0.5,1.03)}, anchor=south, draw=none, font=\scriptsize, /tikz/every even column/.append style={column sep=2pt}},
        legend image code/.code={\draw plot coordinates {(0cm,0cm) (0.3cm,0cm)};}}}
    \subfloat[physical relative to the rule]{%
    \begin{tikzpicture}
        \begin{axis}[scplot, xmin=1, xmax=25, ymin=0.2, ymax=1.5, xtick={1,5,10,15,20,25}, ylabel={ratio [-]}]
            \addplot[reference, forget plot] coordinates {(1,1) (25,1)};
            \addplot[solid, mark=none]          table[x=h, y=Rc] {journals/rhhc/data/scaling_branch.dat};
            \addplot[densely dashed, mark=none] table[x=h, y=Xc] {journals/rhhc/data/scaling_branch.dat};
            \addplot[dotted, line width=0.9pt, mark=none] table[x=h, y=Rl] {journals/rhhc/data/scaling_branch.dat};
            \addplot[dashdotted, mark=none]     table[x=h, y=Xl] {journals/rhhc/data/scaling_branch.dat};
            \legend{$R$ cable, $X$ cable, $R$ line, $X$ line}
        \end{axis}
    \end{tikzpicture}}\hfill
    \subfloat[driving-point impedance at bus~3]{%
    \begin{tikzpicture}
        \begin{semilogyaxis}[scplot, xmin=2, xmax=26, ymin=0.05, ymax=30, xtick={5,10,15,20,25}, ylabel={$|Z_{\Harmonic,33}|$ [pu]}]
            \addplot[gray!60, line width=0.8pt, forget plot] coordinates {(23,0.05) (23,30)};
            \addplot[solid, mark=none]          table[x=h, y=phys]   {journals/rhhc/data/scaling_z.dat};
            \addplot[densely dashed, mark=none] table[x=h, y=scaled] {journals/rhhc/data/scaling_z.dat};
            \legend{physical, scaled}
        \end{semilogyaxis}
    \end{tikzpicture}}
    \caption{Frequency scaling of the fundamental parameters, network model~(iv) at the nominal parameters: (a) series resistance~$R$ and reactance~$X$ of a 500\,mm$^{2}$ cable and a Lion line, physical relative to the scaling rule~$R\sqrt{\Harmonic}$ and~$X\Harmonic$ (gray: rule), and (b) driving-point impedance magnitude at bus~3, physical and scaled (gray: 23rd harmonic, binding limit).}
    \label{fig:scaling}
\end{figure}

\subsection{Computational Tractability}
\label{sec:tractability}

The master problem keeps two variables per harmonic and one cut per scenario, so the effort lies in the oracle. One complete frequency scan of \textsc{LineCableModels.jl} takes about~1.5\,ms per cable design and~2.8--3.6\,ms per overhead-line design, per parameter point and frequency. The full model, with 23 inputs and eight harmonics, requires about 477\,000 such evaluations. With uncertain current angles, every network model converges in two or three iterations to a relative gap below~$5\cdot10^{-4}$, in 34 to 169\,s on 10 threads. The sign-predicted corners reduce the iterations of the full model from five to three and its time by 30\,\%. Without recombination, each earthing variant solves its 49 harmonics in 5.5--7.5\,minutes on 16 threads, in two iterations, to a relative gap below~$10^{-4}$.

\subsection{Comparison with the Literature}
\label{sec:literature_comparison}

With respect to Table~\ref{tab:literature_review}, the robust counterpart solves the 30-bus system in under three minutes with uncertain current angles, and in under eight with 49 harmonics, against under 0.04 minutes for the stochastic problem of~\cite{vanacker_shhc} and 30 minutes for the chance-constrained optimal power flow of~\cite{sun_distribution_2022} on 16 nodes, both with a known network. The studies of the network harmonic impedance~\cite{testa_probabilistic_2002,li_probabilistic_2018,li2020,barakou_parameter_2023} sample or scan the parameters but do not determine a hosting capacity. Whereas an 8\,\% tolerance on the cable capacitance has a minor effect on the harmonic amplification in~\cite{barakou_parameter_2023}, the joint physical inputs cost 14--23\,\% of the hosting capacity here, as the worst cases lie near resonance. The robust methods of~\cite{christakou_voltage_2018,liu_feasible_2023} guarantee feasibility over interval line parameters in linearized models at the fundamental frequency. In this work, the guarantee holds over the exact harmonic network, at the price of an approximate oracle, which the corners and the a-posteriori check control.
\section{Conclusion}
\label{sec:conclusion}

This paper presents a robust harmonic hosting capacity method that combines the uncertain angles of the harmonic currents, treated in probability through chance constraints, with uncertain cable and overhead-line parameters, treated in the worst case over the intervals of their physical inputs. The parameters follow at every harmonic frequency from the physical construction of the cables and lines, and the zero sequence from the windings and earthing of the transformers. With a constant risk factor on exact moments, the robust counterpart is a convex semi-infinite second-order cone program, which a cutting-set method solves with upper and lower bounds. Sign-predicted corners make its worst-case analysis reliable in high dimension, where random sampling misses the worst cases.

On the IEEE 30-bus system with underground cables and overhead lines, the uncertainty of the network costs 14--23\,\% of the hosting capacity up to the 25th harmonic. The cost is uneven over the harmonics, with robust-to-nominal ratios from 0.45 to 0.99, so a robust budget is a new allocation rather than a uniform derating of the nominal one. The worst cases lie at the least damped network, i.e., the coldest cables and lines and the lowest earth resistivity, near the parallel resonances, which the line charging moves and multiplies. Up to the 50th harmonic without recombination, the earthing of the 33\,kV transformers changes the zero-sequence hosting capacity by up to 42\,\% and leaves the positive and negative sequence untouched, and its effect fades at the higher orders, where the zero-sequence current returns through the capacitances of the cables and lines. The Gaussian factor on exact moments stays within 1\,\% of the exact quantile, and its violation probability at the worst cases is at most 1.3 times the risk level. Every case solves within minutes.

These results question the common practice of extrapolating the fundamental-frequency parameters to the harmonics. Scaling the resistance with~$\sqrt{\Harmonic}$ and the reactance with~$\Harmonic$ overestimates the damping of the lines by a factor of two to three and the reactance of the cables by up to 2.4, misplaces the resonance next to the binding limit, and yields an allocation that violates the limits of the physical model, even when made robust over factor intervals derived from the same physical inputs. In the zero sequence, the rule fails altogether, as the resistance follows the return path through the cable screens or the earth rather than the skin effect of the conductor. The harmonic impedance of cables and lines therefore has to be computed from their physical inputs, and the uncertainty of these inputs belongs in the hosting capacity problem itself.

% bibliography
\bibliographystyle{IEEEtran}
\bibliography{journals/rhhc/rhhc_references}

% % biographies
% \begin{IEEEbiographynophoto}{Dr. Ir. Ing. Tom Van Acker} received the M.Eng., M.Sc. and Ph.D. degrees in Electrical Engineering from KU Leuven, Leuven, Belgium in 2012, 2014 and 2020, respectively. Currently, he is a Senior Expert and Team Lead at Elia Transmission Belgium. His main areas of research interest are optimization and uncertainty in a power system context, applied in the context of availability, harmonics, optimal power flow, and state estimation. 
% \end{IEEEbiographynophoto}

% \begin{IEEEbiographynophoto}{Dr.-Ing Amauri G. Martins-Britto} received his BSc, MSc, and PhD in Electrical Engineering from the University of Brasília (UnB) in 2005, 2017, and 2020, respectively. He has been working as a Professional Engineer (P.Eng.) since 2005 and currently he is a Researcher with KU Leuven and the Energy Transmission Competence Hub (Etch) within EnergyVille. His research focuses on power systems modeling, EMT analysis of lines and cables, electromagnetic interferences, and power grounding.
% \end{IEEEbiographynophoto}

% \begin{IEEEbiographynophoto}{Dr. Dipl.-Ing. Hakan Ergun} received his Ph.D. degree in Electrical Engineering for KU Leuven, Belgium in 2015. Currently, he is an Associate Professor with KU Leuven and the Energy Transmission Competence Hub (Etch) within EnergyVille. His research interest is to optimally develop and operate future energy networks for the renewable energy transition. He is the main developer of a number of open-source network modelling tools for optimal planning and operation of AC and DC grids, and stability and security constrained network operation.
% \end{IEEEbiographynophoto}

\end{document}
